% Generated by scripts/paper/supp_experiment_record.py; do not edit.
\paragraph{Pre-SG16 / Pre-SG17 region and second-pass diagnostics (Sec.~94).} Diagnostic only (no method selected). Pre-SG16 asks whether any GT-free region statistic can tell which SPEC32 regions are really foreground, especially regions where the P1 prior is ambiguous (occupancy 0.2-0.8), scored by region ROC-AUC against a GT-majority label. Pre-SG17 asks whether observation agreement predicts whether a pixel change made by the second SA pass is useful (FN$\rightarrow$TP, FP$\rightarrow$TN) or harmful, at fixed threshold 0.45. Run on all 490 VOC / 707 COCO images.

\begin{table}[H]\centering\scriptsize
\begin{tabular}{@{}>{\raggedright\arraybackslash}p{10cm}l@{}}
\toprule
Ambiguous-P1 regions, best occupancy ROC-AUC, VOC & 0.686 \\
Ambiguous-P1 regions, best occupancy ROC-AUC, COCO & 0.673 \\
Ambiguous-P1 regions, best observation statistic ROC-AUC VOC / COCO & 0.597 (obs\_top50\_median) / 0.631 (obs\_all\_median) \\
All regions, p1\_occupancy ROC-AUC VOC / COCO & 0.951 / 0.907 \\
Second-pass useful vs harmful ROC-AUC, second\_q\_direction, VOC / COCO & 0.617 / 0.634 \\
Preregistered SG17 go bar (ROC-AUC) & 0.65 \\
COCO top VOC-frozen reliability decile: useful fraction vs overall & 72.2\% vs 39.0\% \\
Regions evaluated VOC / COCO & 15,680 / 22,624 \\
\bottomrule
\end{tabular}
\end{table}
\emph{Finding.} On ambiguous P1 regions the best simple prior-occupancy ROC-AUC is 0.686 VOC / 0.673 COCO vs 0.597 / 0.631 for the best observation statistic, so no SG16 recovery signal. Best second-pass signal (second\_q\_direction) gives useful/harmful ROC-AUC 0.617 VOC / 0.634 COCO, below the preregistered 0.65 SG17 bar.

\paragraph{Pre-SG18 / Pre-SG21 temporal and head-role diagnostics (Sec.~96).} Diagnostic only. For every one of the 248 call/head observations of SA module 15 (one-pass A\_i c from the class-word seed), and for fixed early/middle/late call stages, measures how many CA-relative true and false positive pixels it adds at fixed threshold 0.45, plus temporal features of each pixel's trajectory and structural descriptors of each head (entropy, local mass, attention distance, reciprocity). Motivated the SG18 window test and the SG21 head weighting.

\begin{table}[H]\centering\scriptsize
\begin{tabular}{@{}>{\raggedright\arraybackslash}p{10cm}l@{}}
\toprule
Stage TP/FP ratio early / middle / late, VOC & 1.354 / 3.494 / 3.228 \\
Stage TP/FP ratio early / middle / late, COCO & 0.509 / 0.941 / 0.935 \\
Added FP $>$16 px fraction early / middle / late, VOC & 0.643 / 0.456 / 0.438 \\
Added FP $>$16 px fraction early / middle / late, COCO & 0.804 / 0.776 / 0.774 \\
Best temporal feature ROC-AUC (middle\_up\_fraction) VOC / COCO & 0.617 / 0.617 \\
VOC-to-COCO rank stability of 248 head/call cells, tp\_fp\_ratio (Spearman rho) & 0.902 \\
Strongest structural correlation, COCO (local\_mass\_r1 vs fp\_separate\_fraction) & 0.949 \\
Images VOC / COCO & 490 / 707 \\
\bottomrule
\end{tabular}
\end{table}
\emph{Finding.} Middle and late stages add far more TP per FP than early calls (VOC TP/FP 1.354 / 3.494 / 3.228; COCO 0.509 / 0.941 / 0.935). Head/call behaviour is highly stable across datasets (Spearman 0.902 for TP/FP ratio over 248 cells), and structural descriptors (local mass, entropy) correlate strongly ($|$rho$|$ about 0.94) with head FP behaviour on COCO. The best temporal pixel feature, middle\_up\_fraction, reaches ROC-AUC only 0.617 on both datasets.

\paragraph{Pre-SG19 late-call support of middle-stage additions (Sec.~99).} Diagnostic only. Takes the pixels that the middle-call map (calls 10-20) adds over the CA seed and asks whether late calls (21-30) can tell the true additions from false ones, which would justify a late-stage veto (SG19). Uses a GPU replay to get all 80 late observations per pixel; feature direction and decile edges are fixed on VOC and transferred to COCO.

\begin{table}[H]\centering\scriptsize
\begin{tabular}{@{}>{\raggedright\arraybackslash}p{10cm}l@{}}
\toprule
Candidate middle-added pixels VOC / COCO & 6,837,692 / 12,671,265 \\
late\_score TP vs FP ROC-AUC VOC / COCO (threshold 0.425) & 0.586 / 0.675 \\
late\_minus\_ca ROC-AUC VOC / COCO & 0.535 / 0.564 \\
late\_up\_consensus ROC-AUC VOC / COCO & 0.556 / 0.530 \\
late\_score decile TP/FP ratio, lowest to highest bin, VOC & 2.71 to 9.57 \\
Same VOC decile edges on COCO, TP/FP ratio & 0.55 to 7.68 \\
VOC person best feature (late\_minus\_middle) ROC-AUC VOC $\rightarrow$ COCO & 0.695 $\rightarrow$ 0.481 \\
Fraction of candidates with L $>$ c, VOC / COCO & 96.8\% / 96.1\% \\
\bottomrule
\end{tabular}
\end{table}
\emph{Finding.} Raw late score separates middle-added TP from FP with ROC-AUC 0.586 VOC / 0.675 COCO, but the specifically temporal features are weak (L - c 0.535 / 0.564) and VOC-person's best signal does not transfer (0.695 $\rightarrow$ 0.481). Verdict: mixed, not a clear go; no SG19 mask was tested.

\paragraph{Pre-SG19 raw-score control (Sec.~100).} Diagnostic only, CPU, from saved SG18 maps. Checks whether the late-score signal in Pre-SG19 is really late-specific, by scoring the same middle-added candidate pixels with five raw scores: the CA seed, the middle map, the late map, and their min and mean.

\begin{table}[H]\centering\scriptsize
\begin{tabular}{@{}>{\raggedright\arraybackslash}p{10cm}l@{}}
\toprule
TP vs FP ROC-AUC, CA c, VOC / COCO & 0.5590 / 0.6198 \\
TP vs FP ROC-AUC, middle M, VOC / COCO & 0.6081 / 0.6861 \\
TP vs FP ROC-AUC, late L, VOC / COCO & 0.5835 / 0.6733 \\
TP vs FP ROC-AUC, min(M,L), VOC / COCO & 0.6006 / 0.6784 \\
TP vs FP ROC-AUC, (M+L)/2, VOC / COCO & 0.5973 / 0.6814 \\
TP vs FP average precision, middle vs late, VOC & 0.8417 vs 0.8262 \\
TP vs FP average precision, middle vs late, COCO & 0.7099 vs 0.6931 \\
Candidate pixels reconstructed VOC / COCO & 6,816,320 / 12,665,100 \\
\bottomrule
\end{tabular}
\end{table}
\emph{Finding.} The middle score beats the late score on both datasets and every contrast (TP/FP ROC-AUC 0.6081 / 0.6861 vs 0.5835 / 0.6733); min and mean do not beat middle. The apparent late signal is explained by general score strength, so the SG19 late veto is not justified.

\paragraph{SG15 region labels vs GT-majority oracle on 20 VOC images (Sec.~93).} Visual diagnostic on 20 VOC-Sim images (the first annotated image of each class, fixed before scoring). For three frozen priors, compares the current SG15 region-relabelling rules with an oracle that labels each SPEC32 region by its GT majority, to see how much headroom hard region decisions have. Not an inference method, no selection, no COCO.

\begin{table}[H]\centering\scriptsize
\begin{tabular}{@{}>{\raggedright\arraybackslash}p{10cm}l@{}}
\toprule
Mean IoU over 20 images, W/P1: prior / A050 / B / oracle (computed) & 0.8031 / 0.7879 / 0.7946 / 0.8690 \\
Mean IoU, W/SG11/trim20/rho0p80: prior / A050 / B / oracle (computed) & 0.7864 / 0.7659 / 0.7838 / 0.8690 \\
Mean IoU, W/SG12/top50: prior / A050 / B / oracle (computed) & 0.7905 / 0.7761 / 0.7875 / 0.8690 \\
Images where oracle IoU $>$ prior IoU, P1 / SG11 / SG12 (computed) & 15 / 17 / 15 of 20 \\
Images & 20 (one per VOC class) \\
\bottomrule
\end{tabular}
\end{table}
\emph{Finding.} Averaged over the 20 images, the GT-majority oracle (0.869 mean IoU) is above every prior and both current rules, while A050 and B both score slightly below the unrefined prior for all three priors. Protocol Section 93 says the oracle values must be presented only as diagnostic upper bounds for hard region decisions on these 20 images.

\paragraph{SG18 middle vs late call-window diagnostic (Sec.~98).} Diagnostic only, CPU, from the cached SG18 maps. Compares one-pass maps from the middle calls (10-20), the late calls (21-30) and both (10-30) at fixed thresholds 0.425 (primary, T1's frozen threshold) and 0.45, and counts which pixels flip when the late calls are added to the middle window. Asks why adding late calls did not help.

\begin{table}[H]\centering\scriptsize
\begin{tabular}{@{}>{\raggedright\arraybackslash}p{10cm}l@{}}
\toprule
mIoU @0.425 middle / late / middle+late, VOC & 0.7959 / 0.7890 / 0.7936 \\
mIoU @0.425 middle / late / middle+late, COCO & 0.6285 / 0.6191 / 0.6240 \\
VOC middle $\rightarrow$ middle+late @0.425: TP gained / lost & 261,001 / 255,053 \\
VOC middle $\rightarrow$ middle+late @0.425: FP gained / removed & 399,322 / 226,040 \\
COCO middle $\rightarrow$ middle+late @0.425: TP gained / lost & 178,332 / 581,249 \\
COCO middle $\rightarrow$ middle+late @0.425: FP gained / removed & 674,271 / 929,652 \\
mIoU @0.45 middle / late / middle+late, VOC & 0.7941 / 0.7868 / 0.7916 \\
mIoU @0.45 middle / late / middle+late, COCO & 0.6343 / 0.6251 / 0.6307 \\
\bottomrule
\end{tabular}
\end{table}
\emph{Finding.} At 0.425, mIoU middle / late / middle+late is 0.7959 / 0.7890 / 0.7936 on VOC and 0.6285 / 0.6191 / 0.6240 on COCO. Adding late calls nets more false positives on VOC but mostly loses true positives on COCO, so protocol Section 98 concludes the effect differs by dataset and cannot be reduced to a single leakage story.

\paragraph{Finalizer reference audit (SG11/SG12 and SG2+DenseCRF recovery) (Sec.~89).} Audit, not an experiment. The SG11/SG12 and SG2+DenseCRF finalizers stopped because new V100 masks did not match historical A100 masks within 0.2\% of pixels at every threshold. This package compares the REF (one-pass) and other control counts across packages and GPU types, and binds the evidence by SHA256, so that within-package V100 comparisons could be recovered without changing any measurement.

\begin{table}[H]\centering\scriptsize
\begin{tabular}{@{}>{\raggedright\arraybackslash}p{10cm}l@{}}
\toprule
Max per-image REF difference, SG11 (V100) vs SG8 (A100), VOC / COCO & 8,185 / 39,917 px \\
Images above the 0.2\% gate, SG11 vs SG8 REF, VOC / COCO & 65 / 168 \\
Max per-image REF difference, two V100 packages (Dense vs SG11), VOC / COCO & 1 / 2 px \\
Max per-image REF difference, A100 SG8 vs SG13, VOC / COCO & 4 / 4 px \\
Max per-image REF difference, SG2+CRF package vs original SG2, VOC / COCO & 26,282 / 18,171 px \\
REF mIoU @0.45, SG11 V100 vs SG8 A100, VOC & 0.788440 vs 0.788656 \\
REF mIoU @0.45, SG11 V100 vs SG8 A100, COCO & 0.625872 vs 0.626588 \\
Live V100 replay (job 234952) of 4 worst samples & SG11/SG12 and raw SG2 counts reproduced exactly \\
\bottomrule
\end{tabular}
\end{table}
\emph{Finding.} Within-package V100 comparisons were recovered; historical pixel-level reproduction was not. Same-hardware runs agree within a few pixels, while cross-hardware REF differences reach 8,185 px (VOC) / 39,917 px (COCO) per image, yet REF mIoU at 0.45 moves by under 0.001. Recovered outcomes (protocol Section 89): SG11 OOF 0.787280 $\rightarrow$ 0.798400, SG12 $\rightarrow$ 0.796658, SG2+CRF does not beat FTF (0.811095 vs 0.813209).

